\section*{使用说明}
本讲义面向已经学过变量、分支、循环、数组和函数，希望系统训练CCF-CSP认证前三题的学生。学习主线是先把短程序写准确，再学习减少重复计算的方法，最后把长题意拆成可靠的程序模块。

全书分成\textbf{第一题、第二题、第三题}三个大的篇章。第一篇依次是“程序基本功与做题技巧、真题速看、考点总结、练习题”；第二篇依次是“基础算法与STL使用、STL、真题速看、考点总结、练习题”；第三篇保留字符串解析、综合模拟案例及综合训练。各篇分别编号，提高专题放入附录。题号是一条学习路线，不表示每次考试的算法固定对应。

\textbf{认证真题}标明认证次数、年月与题号；\textbf{教学例题与自测}用于练习单个知识点。在线做题时注意平台题目的输入格式、数据范围和输出要求。

\textbf{在线练习：}\href{https://oj.shumeng.tech/p}{曙梦OJ题库}；\href{https://oj.shumeng.tech/training/68330d72efce6d5ac2765a62}{历次CCF-CSP认证训练}。可直接打开各算法小节和练习题旁的做题链接。

参考程序以C风格为主，使用scanf/printf、普通数组、下标循环、结构体和独立函数；需要动态容器时再使用STL。STL属于C++，统一按C++17编译。代码有高亮和行号，复制时只包含代码正文（需要阅读器支持）。输入输出样例分区装框，保留空格与换行。

建议先自己推演例题，再说明参考程序的关键不变量，最后遮住答案完成练习。每次写完都应检查最小规模、最大规模、重复值、端点、不可达状态，以及题面允许而样例没有出现的情况。
\tableofcontents
\newpage
